% Autor: Alfredo Sánchez Alberca (email:asalber@ceu.es)
% Gráfico que muestra cómo obtener una muestra a partir de la variable muestral
\begin{tikzpicture}
\tikzstyle{node} = [align=center, node distance=1cm, text=color1]; 
\tikzstyle{arrow} = [-latex, color1, line width=2pt];

\node (x) [label=90:$X$] at (0,5) {\includegraphics[height=1.5cm]{../introduccion/poblacion.png}}; 
\node [anchor=west, align=left] at (6,5) {Población};

\node (x1) [label=90:$X_1$] at (-5,2.5) {\includegraphics[height=1.5cm]{../introduccion/poblacion.png}};
\draw [arrow] (x) -- (x1); 

\node (x2) [label=90:$X_2$] at (-2,2.5) {\includegraphics[height=1.5cm]{../introduccion/poblacion.png}};
\draw [arrow] (x) -- (x2); 

\node [label=90:$n$ copias] at (1,2.5) {\LARGE\ldots};
\node (xn) [label=90:$X_n$] at (4,2.5) {\includegraphics[height=1.5cm]{../introduccion/poblacion.png}};
\draw [arrow] (x) -- (xn); 
\node [anchor=west, align=left] at (6,2.5) {Variable\\ Aleatoria\\ Muestral};

\node (i1) [label=-90:$x_1$, scale=1] at (-5,0) {\includegraphics[height=1.5cm]{../estimacion/individuo1.png}};
\draw [arrow] (x1) -- (i1);
 
\node (i2) [label=-90:$x_2$] at (-2,0) {\includegraphics[height=1.5cm]{../estimacion/individuo2.png}};
\draw [arrow] (x2) -- (i2); 

\node at (1,2.5) {\LARGE\ldots};
\node (in) [label=-90:$x_n$] at (4,0) {\includegraphics[height=1.5cm]{../estimacion/individuo3.png}};
\draw [arrow] (xn) -- (in); 
\node [anchor=west, align=left] at (6,0) {Muestra};
\end{tikzpicture} 